\subsection{APPLICATION: A CRITERION FOR SEMI-SIMPLICITY OF REPRESENTATIONS}

THEOREM 4. Let \(\mathfrak{g}\) be a Lie algebra, \(\mathfrak{r}\) its radical, \(\rho\) a finite-dimensional representation of \(\mathfrak{g},\mathfrak{g}' = \rho (\mathfrak{g})\) and \(\mathfrak{r}' = \rho (\mathfrak{r})\) . Then the following conditions are equivalent:

\begin{enumerate}
\def\labelenumi{(\alph{enumi})}
\item
  \(\rho\) is semi-simple;
\item
  \(\mathfrak{g}'\) is reductive and its centre consists of semi-simple endomorphisms;
\item
  \(\mathfrak{r}'\) consists of semi-simple endomorphisms;
\item
  the restriction of \(\rho\) to \(\mathfrak{r}\) is semi-simple.
\item
  \(\Rightarrow\) (b): if \(\rho\) is semi-simple, \(\mathfrak{g}'\) is reductive (Proposition 5); the associative algebra generated by 1 and \(\mathfrak{g}'\) is semi-simple (Algebra, Chapter VIII, § 5, no. 1, Proposition 3), hence its centre is semi-simple (loc. cit., § 5, no. 4, Proposition 12) and hence the elements of this centre are semi-simple (loc. cit., § 9, no. 1, Proposition 2).
\item
  \(\Rightarrow\) (c): if \(\mathfrak{g}'\) is reductive, its centre is equal to its radical, that is \(\mathfrak{r}'\) , whence the implication (b) \(\Rightarrow\) (c).
\item
  \(\Rightarrow\) (d): suppose that \(\mathfrak{r}'\) consists of semi-simple endomorphisms. As \([\mathfrak{g}',\mathfrak{r}']\) consists of nilpotent endomorphisms (no. 4, Proposition 6), \([\mathfrak{g}',\mathfrak{r}'] = \{0\}\) . Then the implication (c) \(\Rightarrow\) (d) follows from Algebra, Chapter VIII, § 9, no. 2, Theorem 1.
\item
  =\textgreater{} (a): let \(\mathfrak{s}\) be the nilpotent radical of \(\mathfrak{g}\) and \(\rho'\) the restriction of \(\rho\) to \(\mathfrak{r}\) . The elements of \(\rho(\mathfrak{s})\) are nilpotent and hence \(\mathfrak{s}\) is contained in the largest nilpotency ideal of \(\rho'\) . As \(\rho'\) is semi-simple, \(\rho'(\mathfrak{s}) = \{0\}\) and \(\mathfrak{g}'\) is reductive (Corollary to Proposition 6), so that \(\mathfrak{g}' = \mathfrak{a}' \times \mathfrak{r}'\) with \(a'\) semi-simple (Proposition 5). Let \(\mathbf{A}'\) (resp. \(\mathbf{R}'\) ) be the associative algebra generated by 1 and \(a'\) (resp. \(\mathfrak{r}'\) ). It is semi-simple (Algebra, Chapter VIII, § 5, no. 1, Proposition 3), hence \(\mathbf{A}' \otimes_{\mathbb{K}} \mathbf{R}'\) is semi-simple (loc. cit., § 7, no. 6, Corollary 4 to Theorem 3) and hence the associative algebra generated by 1 and \(\mathfrak{g}'\) , which is a quotient of \(\mathbf{A}' \otimes_{\mathbb{K}} \mathbf{R}'\) , is semi-simple, which proves that \(\rho\) is semi-simple.
\end{enumerate}

COROLLARY 1. Let g be a Lie algebra and \(\rho\) and \(\rho'\) two finite-dimensional semi-simple representations of g. Then the tensor product of \(\rho\) and \(\rho'\) is semi-simple.

Let \(\mathfrak{r}\) be the radical of \(\mathfrak{g}\) . For \(x \in \mathfrak{r}\) , \(\rho(x)\) and \(\rho'(x)\) are semi-simple (Theorem 4), hence \(\rho(x) \otimes 1 + 1 \otimes \rho'(x)\) is semi-simple (Algebra, Chapter VIII, § 9, Corollary to Theorem 1) and hence the tensor product of \(\rho\) and \(\rho'\) is semi-simple (Theorem 4).

COROLLARY 2. Let \(\mathfrak{g}\) be a Lie algebra, \(\rho\) a semi-simple representation of \(\mathfrak{g}\) on a finite-dimensional vector space \(V\) , \(T\) and \(S\) the tensor and symmetric algebras of \(V\) and \(\sigma_{T}\) , \(\sigma_{S}\) the representations of \(\mathfrak{g}\) on \(T\) and \(S\) canonically derived from \(\rho\) . Then \(\sigma_{T}\) and \(\sigma_{S}\) are semi-simple and, more precisely, direct sums of finite-dimensional simple representations.

Let \(T^{n}\) be the subspace of T consisting of the homogeneous tensors of order n.~This subspace is stable under \(\sigma_{T}\) and the representation defined by \(\sigma_{T}\) on \(T^{n}\) is semi-simple (Corollary 1). Hence the corollary for \(\sigma_{T}\) and therefore for \(\sigma_{S}\) , which is a quotient representation of \(\sigma_{T}\) .

COROLLARY 3. Let g be a Lie algebra and \(\rho\) and \(\rho'\) two finite-dimensional semi-simple representations of g on spaces M and M'. Then the representation of g on \(\mathcal{L}_{\mathrm{k}}(\mathrm{M},\mathrm{M}')\) canonically derived from \(\rho\) and \(\rho'\) is semi-simple.

The g-module \(\mathcal{L}_{\mathbb{K}}(\mathbf{M},\mathbf{M}^{\prime})\) is canonically identified with the g-module \(\mathbf{M}^{*}\otimes_{\mathbb{K}}\mathbf{M}^{\prime}\) (§ 3, no. 3, Proposition 4), so that Corollary 3 follows from Corollary 1.

COROLLARY 4. Let \(\mathfrak{g}\) be a Lie algebra, a an ideal of \(\mathfrak{g}\) and \(\rho\) a semi-simple representation of \(\mathfrak{g}\) .

\begin{enumerate}
\def\labelenumi{(\alph{enumi})}
\item
  The restriction \(\rho'\) of \(\rho\) to a is semi-simple.
\item
  If \(\rho\) is simple, \(\rho'\) is a sum of simple representations isomorphic to one another.
\end{enumerate}

Passing to the quotient by the kernel of \(\rho\) , \(\rho\) can be assumed to be faithful. Then \(\mathfrak{g}\) is reductive. Let \(\mathfrak{g} = \mathfrak{g}_1 \times \mathfrak{g}_2\) , where \(\mathfrak{g}_1\) is the centre of \(\mathfrak{g}\) and \(\mathfrak{g}_2\) is semi-simple. Then \(\mathfrak{a} = \mathfrak{a}_1 \times \mathfrak{a}_2\) , where \(\mathfrak{a}_1 \subset \mathfrak{g}_1\) , \(\mathfrak{a}_2 \subset \mathfrak{g}_2\) and \(\mathfrak{a}_1\) is the centre of \(\mathfrak{a}\) . The elements of \(\rho(\mathfrak{g}_1)\) , and in particular those of \(\rho(\mathfrak{a}_1)\) , are semi-simple (Theorem 4) and hence \(\rho\) is semi-simple (Theorem 4). Hence (a). Assertion (b) follows from (a), using §3, no. 1, Corollary to Proposition 1.

